\section{Closure}\label{sec:final}

A paper on a problem that remains open has to say where it stops, and saying
so is not a matter of tone. It is a theorem, and this section states it. Every
clause below is proved in the body, none of them depends on a hypothesis that
is not itself proved here or quoted with a reference, and the last clause is
checked by machine.%
\footnote{The quoted results are theorems and are used as such: the
semiregularity theorem of Buchweitz and Flenner \cite{BF08}, the theorem of
Cattani, Deligne and Kaplan on the locus of Hodge classes \cite{CDK95}, the
description of the Hodge rings of abelian varieties of dimension at most five
by Moonen and Zarhin \cite{MZ99}, Markman's theorems on abelian varieties of
Weil type \cite{Mar25a,Mar25b}, and, in the last clause, the theorems of
Deligne, Andr\'e and Milne recorded in \Cref{ssec:bullseye}
\cite{Del71,And96b,Mil20} and Li's theorem on the powers of abelian
fourfolds over finite fields \cite{Li26}, and, for the refutation in
clause (iv), the theorem of Bando and Siu on stable reflexive sheaves
\cite{BS94}, Kuranishi's theorem in the form of Goldman and Millson
\cite{Kur65,GM88}, Grivaux's Chern classes of coherent sheaves \cite{Gri10},
and Voisin's theorem on Weil tori \cite{Voi02b}. \emph{Unconditionally} in \Cref{thm:final}
means that no open statement is assumed anywhere in it; in particular none of
the statements named in its last clause is.} Taken together they fix the boundary of the subject as this
paper leaves it, and the fixing is what we mean by closure: not that the
conjecture is settled, which it is not, but that the distance from here to it
is now measured rather than estimated.

\begin{theorem}[Closure]\label{thm:final}
Let $K=\QQ(\sqrt{-d})$ be imaginary quadratic with $d>0$ squarefree, let
$n\ge2$, and let $\delta\in\QQ^{\times}_{>0}/\Nm_{K/\QQ}(\Kx)$. The following
hold unconditionally.
\begin{enumerate}[label=\textup{(\roman*)},leftmargin=2.8em,itemsep=4pt]
\item \emph{One class.} For an abelian $2n$-fold $A$ of $(K,-1,n)$-Weil type
  with $\Hg(A)=\SU(V,H)$ one has $\dim_{\QQ}\Hdg^{k}(A)=1$ for $k\ne n$ and
  $3$ for $k=n$, and the Hodge conjecture for $A$ in every degree is the
  algebraicity of the single class $\omega_{1}$ of \eqref{eq:omegagens},
  written out with $2^{2n-1}$ integer coefficients
  (\Cref{thm:hdgcount}, \Cref{cor:whatisleft}).
\item \emph{One class, one connected domain, and no loss.} For the family of
  all members of the triple $(K,n,\delta)$ the conjecture is equivalent to the
  algebraicity of $\omega_{s}$ at every point $s$ of the period domain
  $\cD_{n,n}$, which is connected of dimension $n^{2}$; the implication runs
  in both directions, so nothing is given away by passing to this form
  (\Cref{thm:residual}, \Cref{cor:exactform}).
\item \emph{A base point in every family.} Every triple $(K,n,\delta)$
  contains a member whose Weil classes are algebraic, with an explicit
  representative \eqref{eq:divrep} as a rational polynomial in two divisor
  classes, the member being one with quaternionic multiplication extending $K$
  (\Cref{thm:quatweil}, \Cref{thm:everyfamily}). The hypothesis of
  nonemptiness in (ii) is therefore never the obstacle, for any field, any
  dimension or any discriminant.
\item \emph{One finite sufficient condition, false in its first form, and
  one criterion that turns out to be a restatement.} If at $s_{0}$ the
  semiregularity map of a perfect complex whose Chern character is a nonzero
  multiple of a rational Weil class plus a polynomial in the polarisation is
  injective, then $\Sigma=\cD_{n,n}$ and the Hodge conjecture holds in every
  degree for every member of the triple with full Hodge group
  (\Cref{thm:semiregclosure}, \Cref{rem:p2prime}). When the Chern character
  is exactly a multiple of the Weil class, no complex at any member of any
  family has injective semiregularity map (\Cref{thm:p2false}): it would
  deform to very general non-algebraic Weil tori, on which every coherent
  sheaf has vanishing Chern character in degrees $1$ to $2n-1$
  (\Cref{prop:weiltorussheaves}). So the numerical form of that criterion,
  $\dim\Ext^{2}(E,E)=2n(2n-1)$, which would suffice
  (\Cref{thm:p2numerical}) and is the least value possible
  (\Cref{cor:hhfactor}(i)), is attained by no complex, and the shape it
  would force (\Cref{thm:p2endo}) is the shape of nothing. The
  other criterion, that $\omega$ be represented along the orbit by cycles of
  bounded Hilbert data, implies the same conclusion
  (\Cref{thm:degreebound}) and is implied by it
  (\Cref{thm:p1equivalent}), so it is a reformulation of the conjecture for
  the family and not a route to it; read instead on the transported cycle, the
  only representative one can name, it is false whenever a component of the
  base cycle has finite stabiliser (\Cref{thm:p1forces}).
\item \emph{The routes that are closed, and in what sense.} Each of the
  following routes to the Weil class on a general member is closed, in the
  precise sense stated by the theorem cited and delimited in
  \Cref{sec:scope}: a construction is obstructed if it is
  \begin{enumerate}[label=\textup{(\alph*)},nosep,leftmargin=2.2em]
  \item one whose input has Chern character in the subring generated by
    divisor classes, transported by the Poincar\'e kernel
    (\Cref{thm:divisor});
  \item one whose output is a $\Kx$-eigenclass for the norm character
    (\Cref{thm:character});
  \item one carried by an object whose Chern character lies on the Weil line,
    whether through semiregularity or through the weakened criterion, in every
    dimension and by $n^{2}$ dimensions
    (\Cref{thm:rankzero}, \Cref{cor:kernelfamily}, \Cref{thm:weakfails});
  \item one carried by an object with a line bundle summand off the
    polarisation line, hence by any sum of line bundles
    (\Cref{thm:linebundle}, \Cref{cor:noline});
  \item one carried by an object whose Chern character is effective at a base
    point (\Cref{thm:positivity});
  \item induction through hyperplane sections, in either direction
    (\Cref{thm:hyperplane}, \Cref{thm:returntrip});
  \item Kuga--Satake from a weight two sub-structure of $H^{2}(A,\QQ)$, which
    is available only at $n=1$ (\Cref{prop:normpart}, \Cref{cor:ksboundary});
  \item a divisorial argument at a general member of a family of nontrivial
    discriminant (\Cref{prop:divfails});
  \item a secant object whose support fails to be Cohen--Macaulay at a point
    where two components meet only there or where there is an embedded point,
    which removes the candidate of \cite[Section 12]{Mar25b}
    (\Cref{thm:candidatefails}, \Cref{cor:localobstruction},
    \Cref{cor:cmdichotomy});
  \item a secant object of rank one whose support is a complete intersection
    of two divisors, for every $b$ and every $d$ (\Cref{thm:noci});
  \item a secant object quasi-isomorphic to a two-term complex of sums of
    powers of the polarisation, or of simple semi-homogeneous bundles, in
    particular one whose support has a Hilbert--Burch resolution by such
    bundles, in every dimension $n\ge3$ and for every rank, twist and
    discriminant (\Cref{thm:nosplit}, \Cref{thm:noblocks},
    \Cref{cor:disjointroutes});
  \item the semiregularity form of the propagation statement, for any
    representative of the Weil $K$-theory class that is a locally free sheaf
    on a smooth germ of codimension at least two of its support, at every
    member with no abelian subvariety tangent to a $K$-eigenspace; for every
    representative, full rank forces the evaluation map to vanish on the
    $4n^{2}$-dimensional annihilator of the Weil class, which is the mixed
    part of a splitting $HH^{1}(A)=P\oplus Q$ into the annihilators of the two
    conjugate pieces of that class, so that the whole Hochschild action
    factors through a ring of dimension $2^{2n+1}-1$ inside one of dimension
    $2^{4n}$ (\Cref{thm:p2forces}, \Cref{cor:p2shape},
    \Cref{thm:hhsplitting}, \Cref{cor:hhfactor});
  \item the same form of the propagation statement for any representative
    whose support has codimension at least $n$, at every member of every
    family whatever its endomorphisms, hence for every complex carried by a
    cycle that represents the Weil class; and, at a member with no abelian
    subvariety tangent to a $K$-eigenspace, for any representative with a
    component of codimension at least two along which its cohomology has
    nonzero alternating length: two independent jet classes on a complex of
    finite length multiply to a class that the semiregularity map of a torus
    detects through the Euler characteristic, and the Weil line misses every
    class pulled back from a quotient by an abelian subvariety tangent to an
    eigenspace (\Cref{lem:finiteproduct}, \Cref{thm:p2support});
  \item the bounded form of the propagation statement read on the
    transported cycle, for every base cycle with a component of finite
    stabiliser, hence for every cycle this paper writes down: the transport
    multiplies the degree of a component by $c^{2n}/|\ker\varphi\cap G(Z)|$,
    and the denominator is bounded while the numerator is not
    (\Cref{lem:multiplicity}, \Cref{thm:p1forces}); and the same form read on
    a representative chosen afresh at each point, which is equivalent to the
    conjecture for the family and so carries no information
    (\Cref{thm:p1equivalent}).
  \end{enumerate}
\item \emph{Every CM field, not only the quadratic ones.} For every CM field
  $F$ of degree $2m$, every $n\ge2$ and every discriminant, the family of
  abelian varieties of $(F,n,\delta)$-Weil type has a connected period domain
  of dimension $mn^{2}$ carrying a flat space of Weil classes of rank $2m$; it
  contains a member, isogenous to a product of CM abelian $m$-folds, at which
  those classes are polynomials with rational coefficients in divisor classes,
  hence algebraic (\Cref{thm:cmbasepoint}); the algebraic locus is stable
  under the rational points of $\Res_{F_{0}/\QQ}\SU(V,H)$, whose orbits are
  dense, and is a countable union of Zariski closed sets; and each of the two
  criteria of (iv) closes the whole family (\Cref{thm:cmpropagation}). So the
  Weil classes of the higher CM fields are not a separate problem: they are
  the same propagation statement, with the base case supplied
  (\Cref{prop:cmweil}, \Cref{prop:composite}).
\item \emph{What is left, and that it is exactly that.} The Hodge conjecture
  follows from the results of this paper together with three further
  statements, namely the propagation statement of (iv), read for every CM
  field at once, the algebraicity of the Hodge classes
  on abelian varieties that divisor classes and Weil classes do not generate,
  and the conjecture for the varieties that are not abelian. The last of
  these is the conjecture itself, since it covers $A\times\PP^{1}$ for every
  abelian variety $A$ and so contains the abelian case. Its weaker form, the
  conjecture modulo the images of Hodge classes of abelian varieties under
  algebraic correspondences, holds on every variety whose cohomology is
  reached from abelian varieties in that way, in degrees $0$, $2$, $2n-2$ and
  $2n$, and in dimension at most three, and is open beyond. Over the rule set
  of the paper and the literature there are exactly five minimal sets: the
  conjecture for the varieties that are not abelian alone; the Lefschetz
  standard conjecture for every variety with the motivatedness of every Hodge
  class; and the weaker form together with the Lefschetz standard
  conjecture, with the variational Hodge conjecture for algebraic classes, or
  with the propagation statement and the classes beyond the Weil lines, the
  last being the route of this paper. The first two are each equivalent to
  the conjecture, every statement in the five except the propagation
  statement is a consequence of it, and the secant construction of
  \Cref{sec:construction} lies outside all five (\Cref{prop:f3ishc},
  \Cref{def:f3prime}, \Cref{prop:f3prime}, \Cref{thm:closure},
  \Cref{cor:fourremain}, \Cref{prop:lefmot}, \Cref{prop:lefvhc},
  \Cref{thm:vhcab}). Nor are the last two of the three statements
  reductions: the self-product of a Mumford fourfold carries two Hodge classes
  outside the subring of divisor and Weil classes, and the $H^{3}$ of a very
  general quintic threefold is not of abelian type
  (\Cref{prop:mumfordproduct}, \Cref{prop:noabeliantype}). Those two
  classes are the real multiplication by a totally real cubic field on the
  primitive part of $H^{2}$, and they are algebraic once the Kuga--Satake
  class of one K3 surface of Picard number thirteen attached to the fourfold
  is algebraic (\Cref{prop:mumfordrm}, \Cref{thm:mumfordks}). They are
  algebraic at every CM point, and on every member of a Mumford family
  exactly when the Lefschetz standard conjecture holds for the ninefold
  $W\times_{C}W$; one Weil family is likewise equivalent to the Lefschetz
  standard conjecture for one total space (\Cref{thm:mumfordcm},
  \Cref{cor:mumfordequiv}, \Cref{cor:weilequiv}). The Mumford classes are
  rigid, their Hodge locus being the compact base curve, and they are
  algebraic on every member as soon as one perfect complex at one CM point
  has $\Ext^{2}$ of dimension $119$, or as soon as they are represented at
  infinitely many CM points by cycles of bounded degree
  (\Cref{thm:mumfordrigid}, \Cref{thm:mumfordnumerical},
  \Cref{cor:mumfordbounded}); any complex with their Chern character has
  self-extensions of total dimension at least $1600$ when it has no negative
  ones, is not built from line bundles, and extends along no Hochschild
  direction but the curve (\Cref{thm:mumfordprofile},
  \Cref{prop:mumforddivisor}, \Cref{cor:mumfordsplit}). No construction from
  line bundles by cones, tensor products, pull-backs, push-forwards and
  Fourier--Mukai functors reaches them at any point of the curve, and no
  twistor line through a member of the family keeps them of Hodge type; they
  are the twisted dual forms of the part of $H^{2}$ of the square generated by
  its symplectic class, and outside the CM points no rational map from the
  square to a K3 surface or to a hyperk\"ahler eightfold pulls back a
  symplectic form to a nonzero form. At
  every closed point of every reduction of the curve modulo a prime they are
  represented by algebraic cycles, so that they are algebraic on every member
  exactly when those cycles can be chosen with bounded Hilbert data on a
  Zariski dense set of closed points (\Cref{thm:lefschetzclosure},
  \Cref{prop:notwistor}, \Cref{prop:hksquare}, \Cref{prop:modp},
  \Cref{thm:modpdensity}). The routes to them that are not closed off are
  that one statement in three forms and two stronger inputs, any one of
  which would suffice (\Cref{cor:mumfordone}).
\end{enumerate}
\end{theorem}

\begin{proof}
Each clause is the theorem or theorems cited with it, and no clause is used in
the proof of another. For (i), \Cref{thm:hdgcount} counts the Hodge classes and
\Cref{cor:whatisleft} reduces every degree to $\omega_{1}$, the coefficient
count being the rank of $\bigwedge^{2n}H^{1}(A,\ZZ)$ in the basis of
\Cref{thm:explicitgens}. For (ii), \Cref{thm:residual}(a) and (b) are the two
implications and \Cref{prop:weildomain} gives the connectedness and the
dimension. For (iii), \Cref{thm:quatweil}(v) supplies the cycle at a
quaternionic point and \Cref{thm:everyfamily} places such a point in every
family, by the discriminant computation of \Cref{thm:quatweil}(iv) together
with \Cref{lem:products}. For (iv), the sufficient condition is
\Cref{thm:semiregclosure} with \Cref{rem:p2prime}, the refutation of its
first form is \Cref{thm:p2false}, which rests on \Cref{lem:weiltorushodge},
\Cref{lem:weiltorussubsets} and \Cref{prop:weiltorussheaves} and on item
(LI), the lower bound is \Cref{cor:hhfactor}(i), and the sufficiency of the
minimum is \Cref{thm:p2numerical}(iii), with the shape of the object in
\Cref{thm:p2endo}; the bounded
criterion implies the conclusion by \Cref{thm:degreebound}, is implied by it
by \Cref{thm:p1equivalent}, and fails on the transported cycle by
\Cref{thm:p1forces}. For (v), each item is the statement of the theorem
cited, under its own hypotheses, which \Cref{sec:scope} delimits item by
item. For (vi), \Cref{prop:cmweil} gives the family and the reduction to one
class, \Cref{thm:cmbasepoint} the base point and the representative, and
\Cref{thm:cmpropagation} the density and the two criteria. For (vii),
\Cref{prop:f3ishc} shows that the third statement is the conjecture,
\Cref{prop:f3prime} gives what is known of its weaker form,
\Cref{thm:closure}(ii) and (iii) give the five minimal sets and the
necessity of each of their elements, \Cref{thm:closure}(vi) the
consequences of the conjecture, with \Cref{prop:lefmot},
\Cref{prop:lefvhc} and \Cref{thm:vhcab} for the rules of the literature,
and \Cref{thm:closure}(iv) places the secant construction outside them; the computation is item (XXXIII) of
\Cref{sec:verification} and Section 24 of the Lean certificate. The last
two sentences of (vii) are \Cref{prop:mumfordrm} and
\Cref{thm:mumfordks}(ii), with item (XLI), and \Cref{thm:mumfordcm},
\Cref{cor:mumfordequiv} and \Cref{cor:weilequiv}, with item (XLIII); the
added sentence is \Cref{thm:mumfordrigid}, \Cref{thm:mumfordnumerical},
\Cref{cor:mumfordbounded}, \Cref{thm:mumfordprofile},
\Cref{prop:mumforddivisor} and \Cref{cor:mumfordsplit}, with items (XLIV) and
(XLV) and the Lean theorems \texttt{mumford\_rigidity\_sos} and
\texttt{mumford\_ext\_profile}; and the last sentence is
\Cref{thm:lefschetzclosure}, with item (XLVI) and the Lean theorem
\texttt{lefschetz\_counts}, \Cref{prop:notwistor}, with item (XLVII),
\Cref{prop:hksquare}, with item (XLVIII), and \Cref{prop:modp} and
\Cref{thm:modpdensity}, which rest on \cite[Theorem~1.1]{Li26}; the
sentence after it is \Cref{cor:mumfordone}, with item (XLIX).
\end{proof}

\begin{remark}[What this closes]\label{rem:whatclosure}
Four things are settled by \Cref{thm:final}, and further work along the lines
of this paper will not change them.

The problem has a smallest form and it has been reached. By (i) and (ii) the
Hodge conjecture for a whole family of abelian varieties of Weil type is the
algebraicity of one explicitly written class at every point of one connected
domain, and the reduction loses nothing, since the converse holds as well. No
smaller statement is equivalent to it, because the equivalence is already an
equivalence.

The base case is not the difficulty. The treatments of these families in the
literature have a base point only at the split locus, hence only at trivial
discriminant. By (iii) every family has one, for every field, every dimension
and every discriminant, and the cycle is a polynomial in two divisor classes
rather than an existence statement. What separates a base point from the
family is propagation.

Propagation has one criterion, and in its first form it cannot be met. The
bounded criterion is equivalent to its own conclusion, so it restates the
problem rather than reducing it. The semiregularity criterion would be met by
one perfect complex at one base point with $\Ext^{2}$ of dimension
$2n(2n-1)$, but no complex whose Chern character is exactly a Weil class is
semiregular at all, because semiregularity would carry it to non-algebraic
Weil tori, where the Weil class is Hodge and the Chern class of nothing. The
criterion survives only with powers of the polarisation added to the Chern
character, a form this paper does not analyse.

The ways that do not work are known, and they are known for a reason rather
than by exhaustion. Each item of (v) fails for an identifiable structural
cause: a divisor-generated Chern character cannot leave the divisor subring, a
norm eigenclass sits in the wrong summand of the character grading, a Chern
character on the Weil line makes the annihilator too large by exactly the
dimension of the family, a support that is not Cohen--Macaulay carries a
local $\cE xt^{2}$ that the compensating scalar cannot meet, a two-term
complex of powers of the polarisation has summand inclusions that
$H^{2}(\cO)$ does not kill, and a component of nonzero multiplicity makes two
transverse translation classes multiply to a class that the Euler
characteristic detects. These are not obstacles that a more careful version of
the same construction would avoid.

What is not settled is the conjecture, and (vii) says precisely how far from it
this is. Along this route there are three statements, and the third as
first stated, the conjecture for the varieties that are not abelian, is the
conjecture itself (\Cref{prop:f3ishc}). An earlier version of this paper
counted three statements and no fewer; that count rested on a rule set that
omitted the elementary implication of \Cref{prop:f3ishc}, and it was wrong.
With the third statement weakened to the conjecture modulo abelian varieties
the route is minimal again, and three statements remain: the propagation
statement that (v)(l) and (v)(m) constrain and that (vi) carries to every CM
field; the conjecture for the Hodge classes on abelian varieties that no Weil
line reaches, which is not a reduction; and the conjecture modulo abelian
varieties, which holds wherever the cohomology is reached from abelian
varieties by algebraic correspondences and is open beyond
(\Cref{prop:f3prime}, \Cref{rem:f3primeopen}). The shorter routes of the
literature replace the first two by the Lefschetz standard conjecture or the
variational statement, and the Lefschetz standard conjecture is the statement
at their centre (\Cref{rem:bullseye}). The smallest case of the second
statement has been brought to a single correspondence, the Kuga--Satake class
of one K3 surface of Picard number thirteen, and a standard open case of the
third, the square of a K3 surface, meets the same kind of correspondence, the
Kuga--Satake correspondence of that surface; that is the place where a new source of algebraic cycles
would have to act first. We do not
claim any of the three, nor the Lefschetz standard conjecture beyond the
total spaces of Mumford families (\Cref{cor:mumfordlefschetz}), the
motivatedness of Hodge classes or the variational statement, and we have said
so at each place where a reader might reasonably expect otherwise. The value of
(vii) is that the boundary is a computation over a finite rule set rather than
a judgement, so that it can be checked and not merely believed.
\end{remark}
